\chapter{Subversion Command-Line Utilities}\label{svn.utils}

While Subversion\textquotesingle s own \texttt{svn} client is complete and
powerful, a handful of everyday operations---creating and switching
branches, tagging a release, or checking whether a working copy is
clean---take more typing than they ideally would, and their exact
incantations are easy to forget. This appendix describes \emph{Subversion
Utils}, a small, free and open source collection of command-line helpers
that wrap these common operations behind short, memorable commands. They
are written in portable C and run on Linux, macOS, and Windows. The
project lives at
\url{https://github.com/blakemcbride/Subversion-Utils} and is released
into the public domain under the Unlicense.

The utilities assume the conventional Subversion repository
layout---a \texttt{trunk} directory alongside \texttt{branches} and
\texttt{tags} directories, as recommended in
\hyperref[svn.branchmerge]{Branching and Merging}---and they drive the
ordinary \texttt{svn} client under the hood. They therefore require no
special server support and interoperate cleanly with plain \texttt{svn}
usage: you can mix them freely with the standard client, and nothing about
your repository changes by adopting them.

\section{Building and Installing}\label{svn.utils.install}

Building the utilities requires only a C compiler, \texttt{make}, and a
working \texttt{svn} client on your \texttt{PATH}. On Linux and macOS:

\begin{verbatim}
$ make
$ sudo make install            # install into /usr/local/bin
\end{verbatim}

To install elsewhere on your \texttt{PATH} without administrative
privileges, set \texttt{PREFIX}:

\begin{verbatim}
$ make install PREFIX=~/.local
\end{verbatim}

On Windows, build with MinGW-w64:

\begin{verbatim}
$ make CC=x86_64-w64-mingw32-gcc
\end{verbatim}

or, from a Visual Studio Developer Command Prompt, with MSVC:

\begin{verbatim}
> nmake /f Makefile.msvc
\end{verbatim}

Once the utilities are installed, confirm everything is in place by
running \texttt{svn-list} inside a working copy.

\section{The Utilities}\label{svn.utils.commands}

The collection comprises the following commands. Each prints its own usage
message when run with the \texttt{-h} option, even when you are not inside
a working copy.

\begin{description}
\item[\texttt{svn-branch}]
Create, switch, merge, pull, delete, and list branches, and show the name
of the branch the current working copy is on. This collapses the
multistep \texttt{svn\ copy}, \texttt{svn\ switch}, and \texttt{svn\ merge}
dance into a single, easily remembered command.
\item[\texttt{svn-tag}]
Create, switch, delete, diff, and list tags.
\item[\texttt{svn-new}]
Create a fresh repository layout---\texttt{trunk}, \texttt{branches}, and
\texttt{tags}---in a single commit.
\item[\texttt{svn-list}]
Display the repository\textquotesingle s top-level layout (its
\texttt{trunk}, \texttt{branches}, and \texttt{tags}).
\item[\texttt{svn-log}]
Run \texttt{svn\ log} with automatic paging.
\item[\texttt{svn-diff}]
A pager-friendly \texttt{svn\ diff}: it pages its output through
\texttt{less} when run on a terminal, and a single numeric argument is
treated as a revision whose changes should be shown. It can also diff the
\texttt{trunk} or a named branch against your working copy.
\item[\texttt{svn-sync}]
Bring the set of versioned files into line with what is actually present on
disk: schedule unversioned files for addition and missing files for
deletion. (This is unrelated to the \hyperref[svn.ref.svnsync]{svnsync}
repository-mirroring tool.)
\item[\texttt{svn-is-clean}]
Report whether a working copy is clean---that is, free of local
modifications. It is designed for use in scripts and signals its answer
through its exit code: \texttt{0} if the working copy is clean,
\texttt{1} if it contains uncommitted changes, and \texttt{2} if the
target is not a working copy at all.
\end{description}
